% From the WTJ controller to the Special Time Gap controller: design iteration.
% Built by:  uv run iteration.py; uv run final_model.py; uv run build_report.py
% This folder (report/, or report.zip) uploads to Overleaf as is: main.tex, figures/, tables/.
\documentclass[11pt]{article}
\usepackage[margin=1in]{geometry}
\usepackage{amsmath,amssymb}
\usepackage{graphicx}
\usepackage{booktabs}
\usepackage{caption}
\usepackage[hidelinks]{hyperref}

\graphicspath{{figures/final_model/}}
\newcommand{\sat}{\operatorname{sat}}

\title{From WTJ to the Special Time Gap Controller}
\author{}
\date{\today}

\begin{document}
\maketitle

\begin{abstract}
We describe how a smoothing car-following controller for mixed traffic was developed, starting from the
original WTJ controller and ending at a \emph{special time gap} (STG) controller: a constant time gap law
with damping that acts only while the car is closing on the vehicle ahead. Each candidate drove in a
simulated 100-vehicle platoon behind four recorded lead-vehicle drives at four penetration rates, and was
compared with an all-IDM platoon for smoothness (positive work), flow and safety. The iteration showed
that WTJ's speed-dependent target adds nothing at its best setting: the savings come from a time gap law
combined with braking-only damping, a term originally included only to prevent crashes. The final
controller (STG, $\lambda=2$, $\tau=2.5$\,s) cuts positive work by 53\% relative to the all-IDM platoon and
raises flow by 6.6\%. Compared with a plain time gap controller at the same time gap, it uses 17\% less work
for the same flow, and no STG car ever comes within the 5\,m minimum gap.
\end{abstract}

%======================================================================
\section{Setup}
\label{sec:setup}

A leader replays a recorded speed profile, and 100 followers drive behind it in one lane with a 0.05\,s
time step. A share $p$ of the followers (the \emph{penetration rate}) runs the controller under test, and
the rest run the Intelligent Driver Model (IDM: $v_0=30$\,m/s, $T=1.5$\,s, $a_{\max}=0.73$\,m/s$^2$,
$b=1.67$\,m/s$^2$, $s_0=5$\,m). Every controller is run at $p = 0.25$, 0.5, 0.75 and 1 on four drives: the
WTJ real drive (1139\,s), two stop-and-go drives (972\,s and 1227\,s), and a fast mixed drive (1854\,s).
That gives 16 runs per controller. The \emph{baseline} is the all-IDM platoon ($p=0$) on the same drive.
All runs use the same random seed, so the smart cars occupy the same platoon positions for every
controller. The platoon starts at rest behind a parked leader with 50\,m spacing, which causes a strong
start-up transient in the first minute.

\paragraph{Metrics.}
\begin{itemize}
\item \textbf{Work}: the average over followers of $\frac{1}{T}\int_0^T \max(a_i v_i,0)\,dt$, in
  m$^2$/s$^3$. Lower is smoother.
\item \textbf{Flow}: Edie's flow over 250\,m\,$\times$\,30\,s boxes of the time-space plane lying inside
  the platoon, averaged over boxes, in veh/h.
\item \textbf{Throughput}: vehicles per hour passing a fixed point on the road. It is a cross-check only,
  because it rises when the platoon stretches out.
\item \textbf{Safety}: a \emph{collision} ends the run. A run is flagged \emph{$<$5\,m} if a smart car
  comes closer than the minimum gap $g_{\min}=5$\,m.
\end{itemize}
Unless a table states otherwise, results are the percentage change from the baseline, averaged over the
four drives and the four penetration rates.

\paragraph{Notation.}
$s$ is the gap to the car ahead, $v$ the ego speed, and $\dot s = v_{\text{lead}}-v$ the gap rate (negative
while closing). $\sat(\cdot)$ clips the commanded acceleration to $[-3, 1.5]$\,m/s$^2$. Every controller
uses $\alpha=1$\,s$^{-2}$ and $g_{\min}=5$\,m.

%======================================================================
\section{The iteration}
\label{sec:iteration}

Table~\ref{tab:path} lists every candidate in the order it was tried. Only the final controller has no
collisions and no car inside 5\,m, together with a large cut in work and a gain in flow.

\begin{table}[htbp]
\centering
\footnotesize
\caption{Every candidate controller, in the order it was tried. Each was run 16 times (4 drives $\times$ 4 penetration rates). Coll.: runs that ended in a collision. $<$5\,m: runs in which a smart car came closer than the 5\,m minimum gap. Min gap: closest approach of any smart car. Work and flow: change in \%, averaged over drives and penetration rates, against the all-IDM baseline and against TG ($\tau=2.5$\,s) at the same penetration rate; -- if any run crashed.}
\label{tab:path}
\setlength{\tabcolsep}{4pt}
\begin{tabular}{llp{4.6cm}rrrrrrr}
\toprule
 & & & & & & \multicolumn{2}{c}{vs IDM (\%)} & \multicolumn{2}{c}{vs TG (\%)} \\
\cmidrule(lr){7-8}\cmidrule(lr){9-10}
Stage & Controller & What changed & Coll. & $<$5\,m & Min gap (m) & Work & Flow & Work & Flow \\
\midrule
1 & WTJ0 & original WTJ, $\lambda=0.5$ & 16/16 & 0/16 & -- & -- & -- & -- & -- \\
1 & WTJ1 & stronger damping, $\lambda=2$ & 4/16 & 2/16 & 1.15 & -- & -- & -- & -- \\
2 & WTJ2 & soft lower bound, $\beta=0$, $\lambda=3$ & 0/16 & 0/16 & 5.00 & $-$92.3 & $-$37.9 & $-$85.1 & $-$41.9 \\
2 & WTJ3 & $\beta=0.1$ & 4/16 & 0/16 & 5.00 & -- & -- & -- & -- \\
2 & WTJ4 & $\beta=0.25$ & 0/16 & 0/16 & 5.00 & $-$54.9 & $-$5.0 & $-$20.7 & $-$11.4 \\
3 & WTJSoftLB & $\beta=1$ (best of the $\beta$ sweep) & 0/16 & 0/16 & 5.00 & $-$63.0 & $-$2.5 & $-$35.6 & $-$9.0 \\
-- & TG & time gap, two-sided $\lambda=0.1$ (reference) & 0/16 & 12/16 & 0.78 & $-$43.3 & $+$7.0 & ref. & ref. \\
4 & STG & TG with braking-only damping, $\lambda=0.1$ & 0/16 & 12/16 & 0.78 & $-$42.4 & $+$6.8 & $+$1.7 & $-$0.2 \\
\textbf{4--6} & \textbf{STG (final)} & \textbf{braking-only $\lambda=2$, $\tau=2.5$\,s} & \textbf{0/16} & \textbf{0/16} & \textbf{5.00} & \textbf{$-$53.1} & \textbf{$+$6.6} & \textbf{$-$17.3} & \textbf{$-$0.4} \\
\bottomrule
\end{tabular}
\end{table}


\paragraph{Stage 1: original WTJ.}
WTJ follows a gentle speed-dependent target
$a^\star(v) = c\,\big(1/\sqrt{v+0.25} - 1/\sqrt{30.25}\big)$ with $c=0.1$, kept between two time gap laws:
\begin{align*}
a_l &= \sat\!\big(\alpha(s-\tau_{\max}v-g_{\min}) + \lambda\min(0,\dot s)\big), &
a_u &= \sat\!\big(\alpha(s-\tau_{\min}v-g_{\min}) + \lambda\min(0,\dot s)\big), \\
a &= \min\!\big(\max(a_l, a^\star),\, a_u\big), & &
\end{align*}
with $\tau_{\min}=2$\,s and $\tau_{\max}=3$\,s. The work saving was expected to come from following
$a^\star$ inside this band. The braking-only damping $\lambda\min(0,\dot s)$ was added only as a safeguard
against crashes, and was not considered part of the saving. With $\lambda=0.5$ (WTJ0), all 16 runs crashed
during start-up. With $\lambda=2$ (WTJ1), all four runs on the fast mixed drive still crashed mid-drive.

\paragraph{Stage 2: soft lower bound.}
WTJSoftLowerBound pulls the car toward the $\tau_{\max}$ gap instead of holding it there,
$a=\sat\big(\min(\beta a_l+(1-\beta)a^\star,\,a_u)\big)$, and uses $\lambda=3$. At $\beta=0$ (WTJ2) the
cars hardly respond to the car ahead. Work falls by 92\%, but the platoon stretches out and flow falls by
38\%. At $\beta=0.1$ (WTJ3) all four fast mixed runs crashed. $\beta=0.25$ (WTJ4) was the first setting
that was both safe and useful.

\paragraph{Stage 3: $\beta$ sweep.}
$\beta$ was swept from 0 to 1, with $\tau_{\min}=2$ and 2.5\,s. From $\beta=0.2$ upward, work and flow both
improved steadily with $\beta$, so the best setting was $\beta=1$ (WTJSoftLB in Table~\ref{tab:path}). Even so, every setting had 9--44\% less
flow than TG, because the cars are pulled toward a 3\,s gap.

\paragraph{Stage 4: the key observation.}
Because $\tau_{\max}>\tau_{\min}$, $a_l\le a_u$ always. At $\beta=1$, WTJSoftLowerBound therefore reduces to
$a=\sat\big(\alpha(s-\tau_{\max}v-g_{\min})+\lambda\min(0,\dot s)\big)$: a plain time gap law with
braking-only damping. The target $a^\star$ plays no part, so the best WTJ variant got its savings from the
time gap and the damping term. The time gap controller TG used as the reference,
$a=\sat\big(\alpha(s-\tau v-g_{\min})+\lambda\dot s\big)$ with $\lambda=0.1$, damps in both directions.
Making its damping braking-only gives the \textbf{special time gap (STG)} controller:
\begin{equation}
\boxed{a = \sat\!\big(\alpha(s - \tau v - g_{\min}) + \lambda\min(0,\dot s)\big)}
\end{equation}
With TG's small gain ($\lambda=0.1$), STG behaves like TG. With $\lambda=2$, it uses 17\% less work than
TG for the same flow. It also removes TG's close approaches: TG brings a car to 0.78\,m from the car ahead
in 12 of 16 runs.

Table~\ref{tab:compare} shows why STG was chosen over the WTJ variants. WTJ4 and the $\beta=1$ variant
save about as much work or more, but their flow drops as $p$ grows, to 14--16\% below the baseline at
$p=1$. STG has more flow than the baseline up to $p=0.75$ and is within 1\% of it at $p=1$.

\begin{table}[htbp]
\centering
\small
\caption{The final controller (STG, $\lambda=2$, $\tau=2.5$\,s) against the alternatives, by penetration rate $p$: change in work and flow from the all-IDM baseline (\%), averaged over the four drives. TG: time gap controller, $\tau=2.5$\,s. WTJ4: WTJSoftLowerBound, $\beta=0.25$. $\beta=1$: WTJSoftLowerBound, $\beta=1$. Last row: runs (of 16) in which a smart car came closer than 5\,m.}
\label{tab:compare}
\begin{tabular}{lrrrrrrrr}
\toprule
 & \multicolumn{4}{c}{Work (\%)} & \multicolumn{4}{c}{Flow (\%)} \\
\cmidrule(lr){2-5}\cmidrule(lr){6-9}
$p$ & TG & WTJ4 & $\beta=1$ & \textbf{STG} & TG & WTJ4 & $\beta=1$ & \textbf{STG} \\
\midrule
0.25 & $-$12.8 & $-$29.8 & $-$38.8 & \textbf{$-$26.8} & $+$9.0 & $+$6.4 & $+$10.2 & \textbf{$+$13.4} \\
0.5 & $-$38.3 & $-$50.7 & $-$62.3 & \textbf{$-$49.7} & $+$11.3 & $-$0.5 & $+$2.0 & \textbf{$+$10.4} \\
0.75 & $-$54.8 & $-$65.8 & $-$72.9 & \textbf{$-$63.9} & $+$5.8 & $-$9.8 & $-$8.0 & \textbf{$+$3.2} \\
1 & $-$67.4 & $-$73.3 & $-$77.9 & \textbf{$-$72.1} & $+$1.8 & $-$16.3 & $-$14.3 & \textbf{$-$0.6} \\
\midrule
Mean & $-$43.3 & $-$54.9 & $-$63.0 & \textbf{$-$53.1} & $+$7.0 & $-$5.0 & $-$2.5 & \textbf{$+$6.6} \\
\midrule
$<$5\,m & 12/16 & 0/16 & 0/16 & \textbf{0/16} &  &  &  &  \\
\bottomrule
\end{tabular}
\end{table}


\paragraph{Stages 5--6: tuning $\tau$ and $\lambda$.}
Table~\ref{tab:tuning} shows the sweeps. The time gap $\tau$ trades smoothness for flow: between 2 and
3\,s, each extra 0.5\,s cuts work by another 8--11 points and flow by 9--10 points. Where TG runs safely
($\tau\ge2.5$\,s), STG uses 5--10 points less work than TG at the same $\tau$, for about the same flow. STG
is safe at every $\tau$, while TG crashes at $\tau=2$\,s.

The gain $\lambda$ hardly affects flow. In mixed platoons ($p\le0.75$), more damping always saves more
work. In an all-STG platoon ($p=1$), work gets worse again above $\lambda=2$: $-72\%$ at $\lambda=2$, $-69\%$
at $\lambda=3$ and only $-44\%$ at $\lambda=4$. At $\lambda=4$ the cars in an all-STG platoon oscillate.
The likely cause is that a large braking-only gain overshoots, and nothing damps the gap as it reopens.
$\lambda=2$ is the largest gain that does not degrade at $p=1$.

\begin{table}[htbp]
\centering
\small
\caption{Tuning the time gap $\tau$ and the damping gain $\lambda$. Coll.\ and $<$5\,m: runs (of 16) with a collision or with a smart car closer than 5\,m. Work and flow: change from the all-IDM baseline (\%), averaged over drives and penetration rates; Work, $p=1$: the same for an all-smart platoon only. Bold: the final setting.}
\label{tab:tuning}
\begin{tabular}{lrrrrr}
\toprule
Setting & Coll. & $<$5\,m & Work (\%) & Work, $p=1$ (\%) & Flow (\%) \\
\midrule
\multicolumn{6}{l}{\emph{Time gap $\tau$, STG ($\lambda=2$)}} \\
$\tau=2$\,s & 0/16 & 0/16 & $-$42.6 & $-$62.6 & $+$17.0 \\
\textbf{$\tau=2.5$\,s} & \textbf{0/16} & \textbf{0/16} & \textbf{$-$53.1} & \textbf{$-$72.1} & \textbf{$+$6.6} \\
$\tau=3$\,s & 0/16 & 0/16 & $-$60.8 & $-$77.7 & $-$2.3 \\
$\tau=4$\,s & 0/16 & 0/16 & $-$71.0 & $-$83.9 & $-$16.2 \\
\midrule
\multicolumn{6}{l}{\emph{Time gap $\tau$, TG ($\lambda=0.1$, two-sided)}} \\
$\tau=2$\,s & 12/16 & 0/16 & -- & -- & -- \\
$\tau=2.5$\,s & 0/16 & 12/16 & $-$43.3 & $-$67.4 & $+$7.0 \\
$\tau=3$\,s & 0/16 & 0/16 & $-$52.9 & $-$74.5 & $-$2.0 \\
$\tau=4$\,s & 0/16 & 0/16 & $-$66.5 & $-$82.0 & $-$15.6 \\
\midrule
\multicolumn{6}{l}{\emph{Damping $\lambda$, STG ($\tau=2.5$\,s)}} \\
$\lambda=1$ & 0/16 & 0/16 & $-$48.5 & $-$70.0 & $+$6.4 \\
\textbf{$\lambda=2$} & \textbf{0/16} & \textbf{0/16} & \textbf{$-$53.1} & \textbf{$-$72.1} & \textbf{$+$6.6} \\
$\lambda=3$ & 0/16 & 0/16 & $-$55.0 & $-$69.1 & $+$6.4 \\
$\lambda=4$ & 0/16 & 0/16 & $-$50.6 & $-$44.4 & $+$6.2 \\
\bottomrule
\end{tabular}
\end{table}


%======================================================================
\section{The final controller}
\label{sec:final}

The final controller is STG with $\lambda=2$, $\alpha=1$\,s$^{-2}$, $g_{\min}=5$\,m and a selectable time
gap $\tau$, set to 2.5\,s by default. The time gap sets the spacing. The braking-only damping adds the
extra work reduction and the safety margin. $\tau$ sets the trade-off between smoothness and flow.

Table~\ref{tab:final-tau2.5} compares it with TG at the same $\tau$ on every drive and penetration rate.
In every case STG uses 13.5--23\% less work than TG. Its flow is about 4\% higher at $p=0.25$ and up to
5.5\% lower at higher $p$. No STG car ever comes closer than 5\,m, while TG does at every $p\ge0.5$.

\begin{table}[htbp]
\centering
\scriptsize
\caption{The final controller, STG ($\lambda=2$, braking only), against TG at the same time gap $\tau=2.5$\,s, on every drive and penetration rate $p$. Work in m$^2$/s$^3$, flow and throughput in veh/h; $\Delta$: STG relative to TG (\%). The $p=0$ row is the all-IDM baseline. $\dagger$: a TG car came closer than 5\,m (no STG car ever did).}
\label{tab:final-tau2.5}
\begin{tabular}{llrrrrrrrrr}
\toprule
 & & \multicolumn{3}{c}{Work} & \multicolumn{3}{c}{Flow} & \multicolumn{3}{c}{Throughput} \\
\cmidrule(lr){3-5}\cmidrule(lr){6-8}\cmidrule(lr){9-11}
Drive & $p$ & STG & TG & $\Delta$ & STG & TG & $\Delta$ & STG & TG & $\Delta$ \\
\midrule
WTJ real drive & 0 (IDM) & \multicolumn{2}{c}{1.679} & & \multicolumn{2}{c}{1036} & & \multicolumn{2}{c}{702} & \\
 & 0.25 & 1.256 & 1.463 & $-$14.1 & 1171 & 1132 & $+$3.5 & 811 & 779 & $+$4.1 \\
 & 0.5 & 0.843 & 1.021$^\dagger$ & $-$17.4 & 1171 & 1180$^\dagger$ & $-$0.8 & 851 & 853$^\dagger$ & $-$0.2 \\
 & 0.75 & 0.553 & 0.708$^\dagger$ & $-$21.9 & 1109 & 1130$^\dagger$ & $-$1.8 & 838 & 840$^\dagger$ & $-$0.2 \\
 & 1 & 0.396 & 0.470$^\dagger$ & $-$15.9 & 1072 & 1092$^\dagger$ & $-$1.8 & 810 & 813$^\dagger$ & $-$0.4 \\
\midrule
Stop-and-go test (2025-10-29) & 0 (IDM) & \multicolumn{2}{c}{1.317} & & \multicolumn{2}{c}{1053} & & \multicolumn{2}{c}{1060} & \\
 & 0.25 & 0.903 & 1.112 & $-$18.8 & 1191 & 1138 & $+$4.7 & 1397 & 1268 & $+$10.2 \\
 & 0.5 & 0.606 & 0.773$^\dagger$ & $-$21.7 & 1142 & 1161$^\dagger$ & $-$1.6 & 1377 & 1358$^\dagger$ & $+$1.3 \\
 & 0.75 & 0.440 & 0.573$^\dagger$ & $-$23.2 & 1046 & 1106$^\dagger$ & $-$5.5 & 1281 & 1294$^\dagger$ & $-$1.0 \\
 & 1 & 0.323 & 0.376$^\dagger$ & $-$14.0 & 1018 & 1072$^\dagger$ & $-$5.0 & 1189 & 1197$^\dagger$ & $-$0.7 \\
\midrule
Stop-and-go (2026-10-29) & 0 (IDM) & \multicolumn{2}{c}{1.362} & & \multicolumn{2}{c}{1057} & & \multicolumn{2}{c}{1004} & \\
 & 0.25 & 1.011 & 1.201 & $-$15.8 & 1153 & 1108 & $+$4.0 & 1033 & 1014 & $+$1.9 \\
 & 0.5 & 0.691 & 0.849$^\dagger$ & $-$18.6 & 1138 & 1146$^\dagger$ & $-$0.7 & 1094 & 1101$^\dagger$ & $-$0.6 \\
 & 0.75 & 0.475 & 0.595$^\dagger$ & $-$20.2 & 1079 & 1098$^\dagger$ & $-$1.8 & 1075 & 1079$^\dagger$ & $-$0.3 \\
 & 1 & 0.344 & 0.398$^\dagger$ & $-$13.5 & 1042 & 1063$^\dagger$ & $-$1.9 & 1032 & 1039$^\dagger$ & $-$0.6 \\
\midrule
Fast mixed (2026-10-23) & 0 (IDM) & \multicolumn{2}{c}{1.849} & & \multicolumn{2}{c}{1254} & & \multicolumn{2}{c}{1333} & \\
 & 0.25 & 1.388 & 1.650 & $-$15.9 & 1483 & 1430 & $+$3.7 & 1664 & 1655 & $+$0.5 \\
 & 0.5 & 1.002 & 1.200$^\dagger$ & $-$16.5 & 1408 & 1412$^\dagger$ & $-$0.2 & 1519 & 1518$^\dagger$ & $+$0.1 \\
 & 0.75 & 0.798 & 0.953$^\dagger$ & $-$16.3 & 1310 & 1318$^\dagger$ & $-$0.6 & 1422 & 1423$^\dagger$ & $-$0.0 \\
 & 1 & 0.710 & 0.823$^\dagger$ & $-$13.8 & 1238 & 1248$^\dagger$ & $-$0.8 & 1340 & 1340$^\dagger$ & $-$0.0 \\
\bottomrule
\end{tabular}
\end{table}


Figures~\ref{fig:ts}--\ref{fig:work} show the WTJ real drive.
\begin{itemize}
\item \textbf{Time-space diagrams} (Figure~\ref{fig:ts}). The stop-and-go waves of the all-IDM platoon
  fade as the share of STG cars grows. At $p=1$ the platoon passes through each of the leader's slowdowns
  without stopping. TG's diagrams are not shown, because at this scale they look the same as STG's.
\item \textbf{Fundamental diagrams} (Figure~\ref{fig:fd}). At the same density, STG and TG boxes have
  more flow than IDM boxes, and STG and TG overlap.
\item \textbf{Accumulated work} (Figure~\ref{fig:work}). This is where STG and TG differ: at every
  penetration rate, STG's platoon mean ends below TG's. The jump in the first minute is the start-up from
  rest and is the same for both.
\end{itemize}

\begin{figure}[htbp]
\centering
\includegraphics[width=\textwidth]{report_time_space.png}
\caption{Time-space diagrams, WTJ real drive: the all-IDM baseline, then STG ($\lambda=2$,
$\tau=2.5$\,s) at each penetration rate. Color is speed, on one scale for all panels. Thicker lines are
STG cars.}
\label{fig:ts}
\end{figure}

\begin{figure}[htbp]
\centering
\includegraphics[width=\textwidth]{report_fd.png}
\caption{Fundamental diagrams, WTJ real drive, one panel per penetration rate. Each point is one
250\,m\,$\times$\,30\,s box inside the platoon. The grey points (all-IDM baseline) are the same in every
panel.}
\label{fig:fd}
\end{figure}

\begin{figure}[htbp]
\centering
\includegraphics[width=\textwidth]{report_work.png}
\caption{Accumulated positive work per car, WTJ real drive, one panel per penetration rate. Solid lines
are platoon means, and the shaded bands cover the middle 80\% of cars (10th to 90th percentile). The
dashed line is the all-IDM baseline mean.}
\label{fig:work}
\end{figure}

%======================================================================
\section{Limitations}
\begin{itemize}
\item All results come from one random seed, four lead drives and one platoon size. A sweep over seeds is
  needed to put error bars on the differences.
\item The start from rest behind a parked leader is artificial. It causes the WTJ0 and TG ($\tau=2$\,s)
  crashes, and adds the same start-up work to every controller.
\item The explanation of the oscillation for $\lambda\ge3$ is qualitative. A string stability analysis of
  the braking-only damping law would confirm it and give a bound on $\lambda$ as a function of $\tau$.
\end{itemize}

\paragraph{Reproducing the results.}
Run \texttt{uv run iteration.py --results-dir DIR}, then \texttt{uv run final\_model.py --results-dir DIR}
(add \texttt{--tau 2 3} for other time gaps), then \texttt{uv run build\_report.py}. Only missing
simulations are run. \texttt{final\_model.py} also draws the fundamental diagrams and work curves for
every drive and penetration rate in \texttt{plots/final\_model/tau2.5/}.

\end{document}
